% !PRINTED 505
\weberpage{505}

%<<CONT
\gesperrt{by no other operation of $P$ into a determinate quantity $\psi_1$ different from $\psi$.}

Accordingly, to the $j$ subsidiary groups there correspond just as many functions $\psi$, namely
\weq{10}{\psi,\ \psi_1,\ \psi_2\ldots\psi_{j-1},}
and from Theorem 2. there results the theorem:

5. \gesperrt{The quantities (10) undergo a permutation when one and the same operation $\pi$ of $P$ is applied to all of them simultaneously.}

For we can represent the quantities (10) thus
\[
\psi(\varkappa),\ \psi(\varkappa\pi_1),\ \psi(\varkappa\pi_2),\ldots \psi(\varkappa\pi_{j-1}),
\]
and if now $\pi$ is applied to them,
\[
\psi(\varkappa\pi),\ \psi(\varkappa\pi_1\pi),\ \psi(\varkappa\pi_2\pi),\ldots \psi(\varkappa\pi_{j-1}\pi),
\]
in which $\varkappa$ and $\pi$ are arbitrary permutations of $Q$ and $P$. But in each of the two sequences
\[
\begin{aligned}
&\varkappa,\ \varkappa\pi_1,\ \varkappa\pi_2\ldots\varkappa\pi_{j-1}\\
&\varkappa\pi,\ \varkappa\pi_1\pi,\ \varkappa\pi_2\pi\ldots\varkappa\pi_{j-1}\pi
\end{aligned}
\]
there occurs one and only one permutation from each of the subsidiary groups (4).

The functions (10), which, as we have seen, are all different from one another, are called \gesperrt{conjugate functions}.

By what precedes we can easily form the groups to which the individual functions (10) belong. We have, e.g., $\psi_1 = \psi(\pi_1)$, and if now $\pi$ is a permutation of $P$ by which $\psi_1$ remains unaltered, then
\[
\psi(\pi_1\pi) = \psi(\pi_1)
\]
must hold. But we can apply to this equation the permutation $\pi_1^{-1}$ contained in $P$, whereby
\[
\psi(\pi_1\pi\pi_1^{-1}) = \psi
\]
results; that is, $\pi_1\pi\pi_1^{-1}$ must belong to $Q$. If we set it equal to $\varkappa$, it follows that
\weq{11}{\pi = \pi_1^{-1}\varkappa\pi_1.}
If conversely $\pi$ has this form, then
\[
\psi_1(\pi) = \psi(\pi_1\pi) = \psi(\pi_1\pi_1^{-1}\varkappa\pi_1) = \psi(\varkappa\pi_1) = \psi_1,
\]
and $\psi_1$ therefore admits this permutation. Accordingly we can denote the group to which $\psi_1$ belongs by the symbol
\[
\pi_1^{-1}Q\pi_1
\]
%CONT>>
